\section{Related Work}
\label{sec:related}

Our question sits at the intersection of three literatures: capacity scaling laws, mechanistic interpretability of factual knowledge, and probing benchmarks. We summarize each and point out where they fall short of answering ``what fraction of parameters store \emph{standard} knowledge''.

\para{Capacity scaling laws.} \citet{allenzhu2024knowledge} train hundreds of GPT-2-style models on synthetic biographies and derive a tight bit-complexity lower bound, finding that sufficiently-trained transformers store \(\mathbf{2}\) bits of knowledge per parameter --- a ceiling that is universal across GPT-2/LLaMA/Mistral/MoE, preserved at \(\text{int}8\), and halved at \(\text{int}4\). \citet{morris2025memorize} train models on uniform random bitstrings to drive generalization to zero, then read off a raw memorization ceiling of \(\mathbf{3.6}\) bpp. The \(2\) and \(3.6\) bpp figures are not contradictory: the former measures retrievable, structured knowledge tuples; the latter measures raw memorizable bits. \citet{roberts2020packknowledge} provides a real-data counterpart: T5 closed-book QA scales monotonically with parameter count from \(220\)M to \(11\)B. None of these papers partitions the stored bits by \emph{knowledge type}. We use their \(2\) bpp and \(3.6\) bpp ceilings as the denominator in our parameter-share estimate.

\para{Mechanistic interpretability of factual knowledge.} \citet{geva2021ffnkv} establish that the two FFN matrices in each transformer block act as keys (input pattern detectors) and values (output distributions). \citet{dai2022knowledge} apply integrated-gradient attribution and identify roughly four FFN intermediate neurons per fact in BERT-base. \citet{meng2022rome} use causal-mediation analysis to localize factual associations to middle MLP layers at the last subject token, and edit them via a closed-form rank-one update (\rome). \citet{hong2025specialization} show that across 20 open-source LLMs, stronger and more recent models exhibit greater \emph{parameter specialization} --- fewer MLP value vectors per concept, each governing a narrower set of concepts. These results give a per-fact footprint but stop short of an aggregate ``\(X\%\) of parameters'' figure. We invert their per-fact bounds (\(\sim 8\cdot d_{\text{model}}\) parameters per fact for \kn, \(d_{\text{model}}+d_{\text{inter}}\) for \rome) into an analytical upper bound on the share devoted to known facts, which we cross-check against our probe-based estimate.

\para{Probing benchmarks.} \citet{petroni2019lama} introduced \lama, the first cloze-style demonstration that pretrained LMs answer relational queries competitively with traditional KBs. Its limitations --- single-token answers, masked-only formulation, unbalanced answer distribution --- motivated \bear~\citep{wiland2024bear}, which probes by ranking the log-likelihood of full statements rather than predicting a single masked token, and works for both causal and masked LMs without restricting the answer space. \citet{powell2024taxi} introduce \taxi, the first benchmark to evaluate categorical-knowledge consistency under editing: \(976\) edits over \(41\) categories, \(164\) subjects, and \(183\) properties, with \(11{,}120\) MCQ queries. \taxi formalizes ``standard categorical knowledge'' as exactly the property graph our hypothesis discusses (``peanut \(\to\) legume \(\to\) has-shell, contains-protein, edible''). These benchmarks measure \emph{accuracy} but never convert to a parameter footprint. We use \bear-small (\(7{,}731\) instances over \(60\) Wikidata relations) and \taxi baseline-evaluation (\(1{,}435\) categorical queries) and convert per-relation accuracy to acquired bits via the random-baseline-subtraction formula derived in \Secref{sec:method}.

\para{Position of our work.} Unlike capacity-law work, we partition the bit budget by knowledge type. Unlike mechanistic-interpretability work, we report an aggregate parameter share rather than a per-fact footprint. Unlike probe-based work, we convert accuracy into a fraction of model capacity rather than a coverage percentage. The combination yields the missing number: how much of an LLM is standard knowledge?
